\section*{第 \ref{ch:b3:electronic-spectroscopy} 章——电子光谱与光物理}

\begin{solution}{exo:b3:electronic-spectroscopy:1}
(a) 因自旋选择定则禁阻（$\Delta S = 1$）。(b) 因\omterm{thm:b3:electronic-spectroscopy:laporte}{拉波特规则}禁阻
（$g \to g$），但经电子振动耦合可弱见。(c) 允许：$B_{1u}$ 是
$z$ 在 $D_{2h}$ 中的表示。(d) 因对称性禁阻：$A_2$ 不是 $x$、
$y$、$z$ 之一，在 $C_{2v}$ 中；谱带较弱。
\end{solution}

\begin{solution}{exo:b3:electronic-spectroscopy:2}
$10^7/349 - 10^7/450 = 28653 - 22222 = \qty{6431}{cm^{-1}}$, \qty{0.80}{eV}.
\end{solution}

\begin{solution}{exo:b3:electronic-spectroscopy:3}
$A = 5700 \times 1 \times 1.0\times10^{-5} = 0.057$；吸收分数 $1 - 10^{-0.057} =
0.12$.
\end{solution}

\begin{solution}{exo:b3:electronic-spectroscopy:4}
$k_r = \Phi_F/\tau = \qty{1.5e8}{s^{-1}}$, $\tau_r = \qty{6.7}{ns}$; $k_{nr} = (1 - \Phi_F)/\tau
= \qty{1.0e8}{s^{-1}}$.
\end{solution}

\begin{solution}{exo:b3:electronic-spectroscopy:5}
$\int\varepsilon\,\dd\tilde\nu = 1.0645 \times 5700 \times 5000 = \num{3.03e7}$，故 $f \approx
4.32\times10^{-9} \times 3.03\times10^7 = 0.13$：中等强度的
允许跃迁。
\end{solution}

\begin{solution}{exo:b3:electronic-spectroscopy:6}
泊松因子 $\eu^{-S}S^v/v!$ 在 $v = \lfloor S\rfloor$ 处最大。$S = 0.3$ 时，
即 0--0 线本身。$S = 1.5$ 时，$v' = 1$，为 0--0 线的 1.5 倍。$S = 5$ 时，$v' = 4$
与 5 同样大，均为 0--0 线的 26 倍。
\end{solution}

\begin{solution}{exo:b3:electronic-spectroscopy:7}
$\sum k = \qty{4.2e8}{s^{-1}}$: $\tau = \qty{2.4}{ns}$, $\Phi_F = 0.24$, $\Phi_{\mathrm{ISC}} = 0.71$
（且 $\Phi_{\mathrm{IC}} = 0.05$）。
\end{solution}

\begin{solution}{exo:b3:electronic-spectroscopy:8}
$\Phi = 0.546 \times 0.46 \times (0.050/0.040) \times (1.4266/1.333)^2 = 0.36$.
\end{solution}

\begin{solution}{exo:b3:electronic-spectroscopy:9}
$\tau_0/\tau = 1.429$ 与 1.860，对 $[Q]$ 呈线性，故寿命本身
缩短：\omterm{def:b3:electronic-spectroscopy:quencher}{动态猝灭}。$K_{SV} = 43$ L/mol（两点一致），$k_q = 43/8.0
\times 10^{-9} = \qty{5.4e9}{L.mol^{-1}.s^{-1}}$.
\end{solution}

\begin{solution}{exo:b3:electronic-spectroscopy:10}
\omterm{def:b3:electronic-spectroscopy:quencher}{静态猝灭}：仍发射\omterm{def:b3:electronic-spectroscopy:luminescence}{荧光}的分子，寿命与原来相同；其中一半
处于基态，结合于不发\omterm{def:b3:electronic-spectroscopy:luminescence}{荧光}的配合物中。设
缔合常数为 $K_a$，则游离分数为 $1/(1 + K_a[Q])$，且 $I_0/I = 1 +
K_a[Q]$，与\omterm{thm:b3:electronic-spectroscopy:stern-volmer}{斯特恩–福尔默方程}形式相同，但 $\tau$ 不变。
\end{solution}

\begin{solution}{exo:b3:electronic-spectroscopy:11}
$\tau_r = \tau/\Phi_F = \qty{14}{ns}$；$\Phi_T = 1 - 0.36 = 0.64$。实测寿命
因\omterm{def:b3:electronic-spectroscopy:radiationless}{系间窜越}与之竞争而缩短；$\tau_r$ 即仅由发射
所给出的值。
\end{solution}

\begin{solution}{exo:b3:electronic-spectroscopy:12}
$\frac12\langle\alpha\beta - \beta\alpha|\alpha\beta + \beta\alpha\rangle = \frac12(1 + 0 - 0 - 1) = 0$，使用
$\langle\alpha|\alpha\rangle = \langle\beta|\beta\rangle = 1$，$\langle\alpha|\beta\rangle = 0$，对每个电子均成立。
偶极\omterm{def:b3:quantum-model-systems:operator}{算符}不作用于自旋，故\omterm{def:b3:electronic-spectroscopy:transition-dipole}{跃迁偶极矩}中含有此
零因子：$T_1 \leftarrow S_0$ 无强度，除非\omterm{def:b3:many-electron-atoms:spin-orbit}{自旋–轨道耦合}混合了
这些态。
\end{solution}

\begin{solution}{pb:b3:electronic-spectroscopy:1}
\textbf{1.} $A = 5700 \times 1 \times 1.0\times10^{-4} = 0.57$.
\textbf{2.} $1 - 10^{-0.57} = 0.73$.
\textbf{3.} $\varepsilon$ 约为数千，属允许的 $\pi\to\pi^*$ 跃迁，
强度中等。
\textbf{4.} $1239.8/349 = \qty{3.55}{eV}$.
\textbf{5.} 仅吸收紫外线；可见光透过。
\textbf{6.} 为使吸收光与浓度成正比，并避免
发射光的再吸收（内滤效应）。
\textbf{7.} $S_0 \to S_1$（吸收，随后\omterm{def:b3:electronic-spectroscopy:radiationless}{振动弛豫}）；由 $S_1$ 出发：
\omterm{def:b3:electronic-spectroscopy:luminescence}{荧光}、\omterm{def:b3:electronic-spectroscopy:radiationless}{内转换}、\omterm{def:b3:electronic-spectroscopy:radiationless}{系间窜越}至 $T_1$。
\textbf{8.} $k_r = 0.546/19 \times 10^{-9} = \qty{2.87e7}{s^{-1}}$; $\tau_r = \qty{35}{ns}$.
\textbf{9.} $(1 - 0.546)/\tau_0 = \qty{2.39e7}{s^{-1}}$.
\textbf{10.} $28653 - 22222 = \qty{6431}{cm^{-1}}$.
\textbf{11.} 发射始于弛豫后的 $S_1$，终止于
$S_0$ 的振动激发能级，此时分子与溶剂均已弛豫：光子能量较小，
此处足以落入可见光区。
\textbf{12.} $\Phi_F \times (349/450) = 0.42$：占吸收能量的 42\,\%。
\textbf{13.} 数据点每 \qty{0.010}{mol/L} 上升 0.594，是过 1 的直线。
\textbf{14.} 对五个数据点作最小二乘拟合：斜率 $K_{SV} = \qty{59.4}{L/mol}$，
截距为 1.000。
\textbf{15.} $k_q = K_{SV}/\tau_0 = \qty{3.1e9}{L.mol^{-1}.s^{-1}}$.
\textbf{16.} $\tau = \tau_0/(1 + 59.4 \times 0.040) = \qty{5.6}{ns}$.
\textbf{17.} 测量寿命：对于\omterm{def:b3:electronic-spectroscopy:quencher}{动态猝灭}，$\tau_0/\tau$ 落在
与 $I_0/I$ 相同的直线上。
\textbf{18.} 汤力水中不含可使其猝灭的氯离子，且酸性足以
使奎宁保持质子化状态，即其\omterm{def:b3:electronic-spectroscopy:luminescence}{荧光}形式。
\textbf{19.} $k_d = 8 \times 8.314 \times 298/(3 \times 0.890\times10^{-3}) =
\qty{7.4e6}{m^3.mol^{-1}.s^{-1}} = \qty{7.4e9}{L.mol^{-1}.s^{-1}}$.
\textbf{20.} $k_q$ 略小于 $k_d$ 的一半。
\textbf{21.} 约 42\,\%。
\textbf{22.} 更慢：$k_d \propto 1/\eta$，而 $k_q$ 只能随之降低。
\textbf{23.} \textbf{$k_q/k_d = 0.42$}：氯离子猝灭奎宁的速率几乎为
扩散控制碰撞速率的一半。
\end{solution}
